% =============================================================================
%  IEEE Open Journal of the Communications Society (OJ-COMS) framework
%  Compile: pdflatex -> bibtex -> pdflatex x2  (requires IEEEtran.cls)
%  NOTE: this is a FRAMEWORK. Section bodies are stubbed with guidance notes.
% =============================================================================
\documentclass[journal]{IEEEtran}

% --- essential packages ------------------------------------------------------
\usepackage[pdftex]{graphicx}
\usepackage{amsmath,amssymb,amsfonts}
\usepackage{algorithmic}
\usepackage{array}
\usepackage{textcomp}
\usepackage{url}
\usepackage{color}
\usepackage{xcolor}
\def\BibTeX{{\rm B\kern-.05em{\sc i\kern-.025em b}\kern-.08em T\kern-.1667em\lower.7ex\hbox{E}\kern-.125emX}}
\hyphenation{op-tical net-works semi-conduc-tor}

% --- draft todo marker -------------------------------------------------------
\newcommand{\todo}[1]{\textcolor{red}{[TODO: #1]}}

\begin{document}

% =============================================================================
% TITLE BLOCK
% =============================================================================
\title{SINR-Aware Distributed Dynamic Channel Allocation for UAV Swarm
Communications via Behavior Cloning}

\author{
  % Fill in real authors / affiliations per IEEE OJ-COMS author guidelines.
  % OJ-COMS uses the standard IEEE author block (single column, centered).
  Author~One,~\IEEEmembership{Member,~IEEE,}
  Author~Two,~\IEEEmembership{Senior~Member,~IEEE,}
  and~Author~Three,~\IEEEmembership{Fellow,~IEEE}%
  \thanks{Manuscript received ...; revised ...; accepted ....}
  \thanks{A.~One and B.~Two are with the Department of ...,
    (e-mail: ...).}
  \thanks{C.~Three is with ....}
  \thanks{Color versions of one or more figures in this article are available
    online at http://ieeexplore.ieee.org.}
  \thanks{Digital Object Identifier 10.1109/OJCOMS.2026.xxxxxx}
}

% The paper headers (IEEEtran journal mode)
\markboth{Journal of \LaTeX\ Class Files,~Vol.~XX, No.~X, Month~2026}%
{Author \MakeLowercase{\textit{et~al.}}: OJ-COMS Framework}

% =============================================================================
% ABSTRACT + INDEX TERMS
% =============================================================================
\maketitle

\begin{abstract}
\todo{Write a 150--200 word abstract: problem (dynamic distributed channel
allocation for UAV swarms under mutual interference), gap (from-scratch DRL
fails to learn in this high-dimensional, credit-assignment-hard setting),
method (SINR-aware observation + greedy-sequential teacher + behavior cloning
with speed-domain randomization), key results (BC outperforms from-scratch DRL
by up to Xx, matches/beats periodic greedy without needing global CSI,
generalizes zero-shot across scales).}
\end{abstract}

\begin{IEEEkeywords}
UAV swarm communications, distributed channel allocation, behavior cloning,
deep reinforcement learning, interference management, dynamic spectrum access.
\end{IEEEkeywords}

% =============================================================================
% I. INTRODUCTION
% =============================================================================
\section{Introduction}
\label{sec:intro}

\subsection{Motivation}
\todo{Why UAV swarm comms need fast, distributed, interference-aware channel
allocation. Dynamic topology / mobility / mutual interference. Limitations of
centralized control under mobility.}

\subsection{Related Work}
\todo{Centralized allocation; distributed heuristics (graph coloring, JAR);
from-scratch multi-agent DRL (IQL, MADDPG, CARLTON). What is missing: DRL fails
here; no imitation-based solution.}

\subsection{Limitations of Existing Approaches}
\todo{From-scratch DRL: credit assignment + sample inefficiency + exploration
collapse under global interference reward. Heuristics: need global CSI or
suboptimal. Centralized: not deployable on mobile UAVs.}

\subsection{Our Approach and Contributions}
\todo{List 3--4 bullet contributions:
(1) SINR-aware decentralized observation design;
(2) greedy-sequential teacher that provides near-optimal expert actions without
global search;
(3) behavior cloning + speed-domain randomization for dynamic deployment;
(4) extensive comparison showing BC >> from-scratch DRL and competitive with
centralized upper bound, plus zero-shot cross-scale generalization.}

% =============================================================================
% II. SYSTEM MODEL AND PROBLEM FORMULATION
% =============================================================================
\section{System Model and Problem Formulation}
\label{sec:model}

\subsection{UAV Swarm Communication Model}
\todo{Network of N UAVs, shared spectrum, channels, power levels, SINR model.}

\subsection{Interference Model}
\todo{Mutual interference probability, SINR expression, collision definition.}

\subsection{Dynamic Scenario}
\todo{Mobility model, speeds v in {0,10,20,30,40} m/s, time-slotted decisions.}

\subsection{Multi-Agent MDP Formulation}
\todo{State (local SINR/observation), action (channel+power), reward
(interference-aware), decentralized execution assumption.}

\subsection{Problem Formulation}
\todo{Minimize cumulative interference probability P_int subject to
distributed-execution constraint. Define the optimization objective.}

% =============================================================================
% III. PROPOSED METHOD
% =============================================================================
\section{Proposed Method}
\label{sec:method}

\subsection{Framework Overview}
\todo{Two-phase: centralized training (teacher + BC) -> distributed deployment.
Block diagram figure.}

\subsection{Observation Design}
\todo{SINR-aware local observation vector; why it helps credit assignment.}

\subsection{Greedy-Sequential Teacher}
\todo{How the expert generates near-optimal actions sequentially; acts as
upper-bound reference and label source.}

\subsection{BC Training}
\todo{Behavior cloning loss, network, single-variable: speed-domain
randomization (--train_speeds) to remove OOD degradation at high mobility.}

\subsection{Dynamic Deployment}
\todo{Distributed inference at test time, no global CSI, per-slot decisions.}

% =============================================================================
% IV. SIMULATION RESULTS
% =============================================================================
\section{Simulation Results}
\label{sec:results}

\subsection{Experimental Setup}
\todo{Environment params, N in {10,20,30,50}, speeds, baselines (Random,
Dist-Color, Periodic Greedy, CARLTON, IQL, MADDPG, Centralized BC), metrics.}

\subsection{Main Results: Dynamic Scenario}
\todo{Fig.1 + Fig.2: BC vs from-scratch DRL vs heuristics across speeds.}
\subsubsection{Key Observations}
\todo{(1) BC >> from-scratch DRL; (2) MADDPG collapses; (3) periodic greedy
comparable but needs global CSI; (4) BC needs no global CSI.}

\subsection{Scale Comparison}
\todo{Table across N=10/20/30/50; BC stays best, centralized BC upper bound
degrades with N.}

\subsection{Zero-Shot Cross-Scale Generalization}
\todo{BC trained at N=10 generalizes to N=30 (even beats dedicated N=30
centralized BC).}

\subsection{Ablation Studies}
\todo{Effect of speed-domain randomization (s10 vs multi-speed ms); effect of
observation design; teacher quality.}

\subsection{Performance--Computation Trade-off}
\todo{Fig.2 scatter: BC low-latency vs periodic greedy high per-round compute
that scales poorly with N.}

\subsection{Limitations}
\todo{Honest discussion: BC depends on teacher quality; assumptions; failure
modes at extreme scale.}

% =============================================================================
% V. CONCLUSION
% =============================================================================
\section{Conclusion}
\label{sec:conclusion}
\todo{Summarize: imitation (BC) is the right paradigm for this setting where
from-scratch DRL fails; distributed, SINR-aware, mobility-robust; future work
(online refinement, heterogeneous UAVs, real hardware).}

% =============================================================================
% REFERENCES
% =============================================================================
% OJ-COMS uses IEEEtran.bst. Run: pdflatex -> bibtex -> pdflatex x2
\bibliographystyle{IEEEtran}
\bibliography{paper_ojcoms}

\end{document}
